\scanpage{32}
\mspage{19}

partie ouverte $I'$ de $I$. \deleted{\unsure{Si}}

On prouvera en $=$ \deleted{\unsure{temps}} le lemme et son corollaire, en procédant par récurrence sur $\mathrm{card}.\ I=n$. Si $n\unsure{\leq|=}0$ i.e. $I=\varnothing$, les énoncés sont triviaux. Par transitivité, en utilisant les diagrammes du type \unsure{(1?)|(1)?|(2?)}, on est ramené au cas où $I=I'\cup\{i_0\}$, où $i_0\in I\smallsetminus I'$. On pose \inserted{encore} $I_0=I_{/\unsure{i|l}_0}$, $I'_0=I_0\cap I'$, et le raisonnement qu’on vient de faire nous montre qu’il suffit de prouver que $u_{I'_0,I_0}\in C$ resp. $\in TC$. Si $I'_0\ne I_{\deleted{\unsure{0|g}}}$, cela résulte de l’hyp. de récurrence. Sinon, $i_0$ est un objet final de $I$, et $I'=I_{<i_0}$, et alors $u_{I',I}\in C$ \deleted{résulte} est inclus dans la définition de $C_I$. \deleted{\unsure{Soit}} Reste à prouver \inserted{donc} que $u_{I',I}\in TC$ si les $u_i\in W$. On a le diagramme

\diagram{\begin{tikzcd}[ampersand replacement=\&,column sep=large]
 \unsure{X}_{I'}\arrow[r]\arrow[d,"u_{I'}"']\&X_I\arrow[d,"{\widetilde u_{I',I}}"]\arrow[dd,bend left=45,dashed,"u_I"]\\
 Y_{I'}\arrow[r]\&Y_{I',I}\arrow[d,"{u_{I',I}}"]\\
 \&Y_I\ ,
\end{tikzcd}}
% archive sheet number: 31

